\documentclass[11pt]{article}

\usepackage[margin=1in]{geometry}
\usepackage[T1]{fontenc}
\usepackage[utf8]{inputenc}
\usepackage{amsmath,amssymb}
\usepackage{graphicx}
\usepackage{booktabs}
\usepackage[round]{natbib}
\usepackage{hyperref}
\hypersetup{colorlinks=true, linkcolor=blue!50!black, citecolor=blue!50!black, urlcolor=blue!50!black}
\usepackage{enumitem}
\usepackage{xcolor}

\title{\LARGE Prompt-Driven Interactive Segmentation\\[4pt]
\large as a Dataset-Annotation Tool}
\author{
  \begin{tabular}{c}
    Student A: 2301498\\
    Abdelhafeez Aref Abdelhafeez Almaz
  \end{tabular}
  \and
  \begin{tabular}{c}
    Student B: 2501925\\
    Abdelrahman Eid Ashour Mohamed
  \end{tabular}
}
\date{Milestone 1 --- October 31, 2026\\CSE 681 --- Digital Image Processing}

\begin{document}
\maketitle

\begin{abstract}
Producing accurate instance masks for image datasets is time-consuming when
every object must be outlined manually. We propose a local interactive
annotation tool that uses positive/negative point and box prompts with the
frozen SAM 2.1 Hiera-Tiny model, and measures annotation time and mask agreement
with existing ground truth. We will use a deterministic 300-image subset of
COCO 2017 val2017, retain its non-crowd instance masks, and compare prompted
annotation against manual polygon labeling on the same target objects. The
project will report per-instance intersection over union (IoU), Dice, annotation
time, and inference latency without training a model or claiming results before
they are measured.
\end{abstract}

\section{Problem Statement \& Motivation}
\label{sec:motivation}
Instance segmentation assigns a pixel mask to each object in an image, but
creating these masks by hand is labor-intensive and can limit the scale of
downstream image-processing experiments. This project investigates whether
prompt-driven segmentation can reduce human annotation effort while retaining
useful agreement with reference masks. The work connects directly to the course
topics of digital image segmentation, region delineation, binary-mask
representation, and quantitative evaluation. Rather than train a new
segmentation model, we will build a reproducible human-in-the-loop tool around a
frozen pretrained model and evaluate both the resulting masks and the time
required to accept them.

\section{Related Work}
\label{sec:related}
Promptable segmentation builds on interactive segmentation and recent
foundation-model approaches. GrabCut established a classic interactive
foreground-extraction method in which user input guides iterative graph-cut
optimization \citep{rother2004grabcut}. The Segment Anything Model (SAM)
introduced a promptable model trained at large scale and evaluated with
point/box prompts \citep{kirillov2023sam}. SAM 2 extends promptable segmentation
to images and videos with a streaming memory architecture; we use its image
predictor rather than training a model from scratch \citep{ravi2024sam2}.
FastSAM is a recent alternative that seeks faster segment-anything behavior
with a different model pipeline \citep{zhao2023fastsam}.

Our project is a small systems and evaluation study, not a new foundation model
or a reproduction of SAM 2 training. COCO provides the images and instance masks
needed for a post-hoc comparison \citep{lin2014coco}. We will compare an
interactive SAM 2.1 workflow with manual polygon annotation on matching target
objects. The evaluation will separate model inference latency from end-to-end
annotation time, and will use the held-out COCO mask only for scoring, not to
choose prompts or candidate masks.

\begin{table}[ht]
\centering
\caption{Key references.}
\label{tab:refs}
\begin{tabular}{@{}p{0.29\linewidth}p{0.08\linewidth}p{0.53\linewidth}@{}}
\toprule
\textbf{Reference} & \textbf{Year} & \textbf{Relevance to this project} \\
\midrule
\citep{rother2004grabcut} & 2004 & Classic interactive foreground extraction guided by user input. \\
\citep{lin2014coco} & 2014 & Public images and instance-segmentation ground truth for evaluation. \\
\citep{kirillov2023sam} & 2023 & Promptable segmentation with points, boxes, and masks. \\
\citep{zhao2023fastsam} & 2023 & Recent promptable segmentation alternative for related-work context. \\
\citep{ravi2024sam2} & 2024 & Frozen image/video segmentation model used by our tool. \\
\bottomrule
\end{tabular}
\end{table}

\section{Dataset \& Baseline Plan}
\label{sec:data}
We will sample 300 images from COCO 2017 val2017. Candidate images must have
between one and eight non-crowd instance annotations; candidates are sorted and
sampled with Python's pseudorandom generator using seed 42. We will record the
selected image IDs, retain the corresponding instance annotations, and exclude
crowd regions from scoring. The COCO images retain their individual source
licenses, so we will follow the license metadata for every image and will not
redistribute the images or derived masks. COCO's terms are documented by the
dataset maintainers \citep{cocoTerms}.

The primary baseline is manual polygon annotation of the same target object in
each evaluated image, using a standard polygon annotation interface without SAM
assistance. Student A will record active annotation time using the same start
and stop definition used by the interactive tool. For the prompted method,
Student B will record time from the first prompt through accepting the mask;
inference latency will be logged separately. Before either method is run, the
target COCO instance ID will be recorded. After each mask is finalized, we will
compare it with its matching COCO non-crowd instance mask. The reference mask
will not be exposed while prompts are chosen.

\begin{table}[ht]
\centering
\caption{Dataset summary.}
\label{tab:data}
\begin{tabular}{@{}p{0.20\linewidth}p{0.43\linewidth}p{0.27\linewidth}@{}}
\toprule
\textbf{Dataset} & \textbf{Subset used} & \textbf{Link / terms} \\
\midrule
COCO 2017 & 300 val2017 images; seed 42; 1--8 non-crowd instances per image & \url{https://cocodataset.org/} \\
\bottomrule
\end{tabular}
\end{table}

\section{Proposed Method}
\label{sec:method}
The application has a React/TypeScript browser interface and a Python/FastAPI
backend. The backend runs the frozen SAM 2.1 Hiera-Tiny image predictor with
positive or negative point prompts and/or a box prompt. For each request, the
predictor returns candidate masks and confidence scores; the tool displays the
highest-scoring candidate for human review. The annotator can refine prompts,
accept the mask, and save it. The evaluation pipeline associates the accepted
mask with its COCO annotation ID, then computes binary-mask IoU and Dice against
the corresponding ground truth. The annotation timer starts at the first prompt
and ends when the annotator accepts the mask; inference time is measured
separately. The main contribution is a small, reproducible
annotation-and-measurement workflow, not a change to SAM 2's model architecture.

\section{Feasibility \& Compute Plan}
\label{sec:feasibility}
The project is CPU-OK because it performs inference with a frozen, small SAM 2.1
checkpoint and does not train or fine-tune the model. The selected Hiera-Tiny
checkpoint is approximately 149 MiB; inference latency will be measured on the
available machine and reported rather than assumed. The 300-image subset is
small enough for local storage and evaluation. If CPU inference is too slow for
the full study, we will first complete the MVP on a fixed smaller sample and
report the actual sample size and device.

\subsection{Minimum Viable Deliverable (MVP)}
Run the interactive tool end-to-end on a fixed sample of 30 COCO images; measure
manual polygon and prompted annotation time on matching target objects; save
both methods' masks; and report mean/median annotation time, mean IoU and Dice,
and per-image results. The sample IDs, prompts, device, model checkpoint, and
timing protocol will be recorded so the comparison can be reproduced.

\subsection{Stretch Goal}
Expand the comparison to all 300 images, report results by object category and
prompt type (click versus box), include representative failure cases, and run a
small usability test with a second annotator.

\section{Risk Assessment}
\label{sec:risk}
\begin{itemize}[leftmargin=*]
  \item \textbf{CPU latency or memory limits:} use the Hiera-Tiny checkpoint,
  measure inference separately from annotation time, and reduce the evaluated
  sample to the predeclared MVP if needed; no training is required.
  \item \textbf{Dataset retrieval or license variation:} download only the
  selected COCO images from the official source, preserve the selected IDs and
  annotation metadata, and follow each image's source license; do not commit or
  publicly redistribute dataset assets.
  \item \textbf{Prompt and ground-truth mismatch:} identify the target instance
  before annotation, keep its COCO annotation ID with the result, and score only
  after acceptance; crowd regions and empty masks are excluded.
\end{itemize}

\section{Work Breakdown}
\label{sec:workbreakdown}
\begin{table}[ht]
\centering
\caption{Milestone task allocation. Student B combines the catalog's model and evaluation roles.}
\label{tab:tasks}
\begin{tabular}{@{}p{0.12\linewidth}p{0.39\linewidth}p{0.39\linewidth}@{}}
\toprule
\textbf{Milestone} & \textbf{Student A --- Data \& Baseline} & \textbf{Student B --- Model + Evaluation} \\
\midrule
M1 (Wk 3) & Prepare the COCO subset and ground-truth masks; define the manual-polygon timing protocol. & Build the click/box interface, run SAM 2 inference, and define the IoU/time evaluation protocol. \\
M2 (Wk 6) & Measure manual polygon time on the shared sample and verify target-mask IDs. & Deliver the interactive tool and first IoU/time comparison; combine catalog roles B and C. \\
M3 (Wk 9) & Complete manual-baseline timing and data checks. & Polish the tool; complete the full comparison, failure analysis, and small user test. \\
\bottomrule
\end{tabular}
\end{table}

\section*{AI Tool Usage Disclosure}
AI assistance was used for initial code scaffolding and drafting/editing this
proposal. The team will verify the implementation, citations, measurements, and
final technical claims, and will revise the report to reflect work they have
personally reviewed and understand.

\bibliographystyle{unsrtnat}
\bibliography{references}

\end{document}
